\documentclass[11pt,a4paper,titlepage]{article}
\usepackage[margin=2.5cm]{geometry}
\usepackage[style=bath, backend=biber]{biblatex}
\addbibresource{references.bib}

\title{Working Title Goes Here \\[1em] \large Project Proposal}
\author{Matthew La \\ Bachelor of Science in Computer Science \\ The University of Bath}
\date{\today}

\begin{document}

\pagenumbering{roman}
\maketitle
\tableofcontents

\newpage
\pagenumbering{arabic}

\section{Problem Description}

Since the public availability of GitHub Copilot and ChatGPT in 2022, AI-assisted code generation has become a staple of software engineering, signalling a shift away from hand-written code. In 2026, Stack Overflow's Developer Survey reported that 83\% of respondents used AI tools at work \parencite{stackoverflow2026}. Furthermore, GitHub also reports that the proportion of pull requests that are initially drafted by AI has more than doubled between March and September 2026 \parencite{github2026octoverse}.

However, studies have shown that much of AI-generated code can be incorrect, insecure, or both. For example, \textcite{vero2025baxbench} show that even OpenAI o1 may perform poorly when writing realistic software, with only 62\% code correctness. Additionally, roughly half of these correct programs still contained vulnerabilities.

Trust in LLM output therefore remains a significant factor for AI adoption. Indeed, while 87\% of Stack Overflow survey respondents report some level of trust in AI's outputs, 48\% trust AI only when they can easily verify its output. 

One of the most common methods for developers to carry out this kind of verification for their software is unit testing. Even so, test-passing AI-generated software patches still contain errors at a noticeable rate, with 29.6\% of these patches behaving differently from the developers' real fix, and 28.6\% of these divergences containing incorrect code \parencite{wang2025solvedissues}.

More fundamentally, tests are a finite set of hand-picked inputs that inevitably miss unconsidered edge cases. As \textcite[p.~7]{dijkstra1970notes} observed, ``program testing can be used to show the presence of bugs, but never to show their absence''. Dijkstra's own later solution to this problem was to explicitly prove a program's correctness \parencite{dijkstra1972humble}.

\textcite{hoare1969axiomatic} gave this idea a formal basis, and formal verification can now guarantee that a program meets its specification for every possible input. However, these guarantees are expensive. The seL4 microkernel's roughly 8,700 lines of C required around 200,000 lines of Isabelle/HOL proof \parencite{klein2009sel4}, and a survey of 216 practitioners identifies scalability and the expertise required as key barriers to adoption \parencite{gleirscher2020formal}.

Yet the idea at the core of verification, stating explicitly what a program should do and checking the code against that statement mechanically, need not carry this cost. Design by contract \parencite{meyer1992dbc} attaches such specifications to each function as preconditions and postconditions, making intended behaviour part of the code rather than leaving it implicit in a set of tests. AI-generated code may make this more valuable still \parencite{kleppmann2025verification}: when an LLM writes the implementation, a contract records intent independently of it, giving the compiler something to check the code against.

\subsection{Aims}
\subsection{Objectives}

\section{Requirements Specification}

\section{Project Plan}
\subsection{Gantt Chart}
\subsection{Tasks and Milestones}
\subsection{Risk Assessment}

\section{Resources}

\newpage
\printbibliography

\end{document}
